\centerline{\bf \Huge Домашнее задание.}
\centerline{\bf \Huge Лемма о венике}

\begin{flushright}
        \scriptsize{А ведь вы могли выбрать антипараллельность и быть счастливыми...}
\end{flushright}


\problem{1!}
Центр вписанной окружности треугольника $ZOV$ находится в точке $I$.
К описанной окружности треугольника $OIV$ проведенная касательная, проходящая через точку $I$ и пересекающая стороны $ZO$ и $ZV$ в точках $K$ и $L$. соответственно. Докажите, что $ZK = ZL$.

\problem{2!}
Дана равнобокая трапеция $GOAL$ с основаниями $OA$ и $GL$. Точки $I_1, I_2$ являются центрами вписанных окружностей треугольников $GOA$ и $GOL$ соответственно. Докажите, что $I_1I_2 \perp OA.$

\problem{3!}
Точки $A,B$ и $C$   лежат на окружности с центром в точке $O$. Луч $OB$   вторично пересекает описанную около треугольника $AOC$ окружность в точке $D$,   причём точка B   оказалась внутри этой окружности. Докажите, что $AB$ --- биссектриса угла $DAC$.

\problem{4}
В треугольнике $ABC$ $( AB < BC)$  точка $I$ --- центр вписанной окружности, M – середина стороны $AC, N$ --- середина дуги $ABC$ описанной окружности.
Докажите, что  $\angle IMA = \angle INB$.

\problem{5}
Отрезки $AB$ и $CD$ длины $1$ пересекаются в точке $O$, причём $\angle AOC = 60^{\circ}$. Докажите, что $AC + BD > 1$.

\problem{6}
$\bigl[\textbf{ЕГЭшечка}\bigr]$
Окружность, вписанная в треугольник $MNK$, касается сторон $MN, NK$ и $MK$ в точках $A, B$ и $C$ соответственно.

а) Докажите, что $NB = \dfrac{MN+NK-MK}{2}$

б) Найдите отношение $MA:AN$, если известно, что $NB:NK = 1:3$ и $\angle MNK = 60^{\circ}$